%% Source: MANUSCRIPT.md "Materials and Methods > Atom mapping"
%%         (drafts/methods_atom_mapping.md).
\subsection{Atom mapping}\label{sec:methods-atom-mapping}

Atom mapping --- the assignment of individual reactant atoms to their
product-side counterparts across a mass-balanced reaction --- is a new
capability in the 2026 release. Mappings enable atom-level tracing of carbon,
phosphorus, sulfur, and other atoms through a metabolic network, and are
increasingly required for downstream applications including
\textsuperscript{13}C metabolic flux analysis, pathway-degradation prediction,
and cofactor-usage attribution.

Mappings for the ModelSEED biochemistry are being generated through a
collaboration with the \TBD[group] group at the \TBD[institution]. Mapping
methodology: \TBD[will be described briefly with reference to the
collaborators' methods paper]. Delivery format: \TBD[likely RXN files or a
tabular format per reaction].

Mapped reactions are integrated into ModelSEED under
\file{Biochemistry/AtomMappings/} (or equivalent, \TBD{}). Each mapped reaction
is stored alongside its ModelSEED reaction ID; downstream consumers can access
mappings through the same Solr endpoint or the new PyPi API (see
Section~\ref{sec:methods-distribution}).

Coverage as of the 2026 snapshot: \TBD[fraction of priority-scope reactions with
mappings delivered]. Coverage of the full ModelSEED reaction set: \TBD{}.
Reactions without a mapping fall into three classes: (i) reactions whose
reactant compounds do not yet have a stored structure (necessarily unmappable);
(ii) reactions with \texttt{R} groups whose mapping requires expansion to
specific chain lengths; (iii) reactions whose complexity (e.g. very long acyl
chains, PKS/NRPS iterative modules) exceeds the mapper's practical scope.
